\Section{Number Theory}

\Topic{Linear diophantine equation}. $ax+by=c$.
Let $d=\gcd(a,b)$. A solution exists iff $d|c$.
If $(x_0,y_0)$ is any solution, then all solutions are given by
$(x,y) = (x_0 + \frac{b}{d}t, y_0 - \frac{a}{d}t)$, $t \in {\mathbb Z}$.
To find some solution $(x_0, y_0)$, use extended GCD to solve
$ax_0 + by_0 = d = \gcd(a, b)$, and multiply its solutions by $\frac{c}{d}$.

Linear diophantine equation in $n$ variables:
$a_1 x_1 + \dots + a_n x_n = c$ has solutions iff $\gcd(a_1, \dots, a_n) | c$.
To find some solution, let $b=\gcd(a_2, \dots, a_n)$,
solve $a_1 x_1 + by = c$, and iterate with $a_2 x_2 + \dots = y$.

\Topic{Extended GCD}
\begin{verbatim}
// Finds g = gcd(a,b) and x, y such that ax+by=g.
// Bounds: |x|<=b+1, |y|<=a+1.
void gcdext(int &g, int &x, int &y, int a, int b)
{ if (b == 0) { g = a; x = 1; y = 0; }
  else        { gcdext(g, y, x, b, a % b); y = y - (a / b) * x; } }
\end{verbatim}

Multiplicative inverse of $a$ modulo $m$:
$x$ in $ax + my = 1$, or $a^{\phi(m)-1} \pmod{m}$.

\Topic{Chinese Remainder Theorem}.
System $x \equiv a_i \pmod{m_i}$ for $i=1,\dots,n$, with
pairwise relatively-prime $m_i$ has a unique solution modulo $M = m_1 m_2 \dots m_n$:
$x = a_1 b_1 \frac{M}{m_1} + \dots + a_n b_n \frac{M}{m_n} \pmod{M}$,
where $b_i$ is modular inverse of $\frac{M}{m_i}$ modulo $m_i$.

%\Topic{Generalized CRT}.
System $x \equiv a \pmod{m}$, $x \equiv b \pmod{n}$ has solutions
iff $a \equiv b \pmod{g}$, where $g=\gcd(m,n)$.
The solution is unique modulo $L=\frac{mn}{g}$, and equals:
$x \equiv a + T(b-a) m/g \equiv b + S(a-b) n/g \pmod{L}$,
where $S$ and $T$ are integer solutions of $mT + nS = \gcd(m,n)$.

\Topic{Prime-counting function}. $\pi(n) = |\{p \le n : p \hbox{ is prime}\}|$.
$n/\ln(n) < \pi(n) < 1.3 n/\ln(n)$.
$\pi(1000) = 168$, $\pi(10^6) = 78498$, $\pi(10^9) = 50\ 847\ 534$.
\quad $n$-th prime $\approx n \ln n$.

\Topic{Miller-Rabin's primality test}.
Given $n = 2^r s + 1$ with odd $s$, and a random integer $1 < a < n$. \\
If $a^s \equiv 1 {\pmod n}$ or $a^{2^j s} \equiv -1 {\pmod n}$ for some
$0 \le j \le r-1$, then $n$ is a probable prime.
With bases 2, 7 and 61, the test indentifies all composites below $2^{32}$.
Probability of failure for a random $a$ is at most 1/4.

\Topic{Pollard-$\rho$}.
Choose random $x_1$, and let $x_{i+1} = x_i^2 - 1 \pmod{n}$.
Test $\gcd(n, x_{2^k+i} - x_{2^k})$ as possible $n$'s factors for $k=0,1,\ldots$
Expected time to find a factor: $O(\sqrt{m})$, where $m$ is smallest
prime power in $n$'s factorization.
That's $O(n^{1/4})$ if you check $n = p^k$ as a special case before factorization.

\Topic{Fermat primes}.  A Fermat prime is a prime of form $2^{2^n}+1$.
The only known Fermat primes are 3, 5, 17, 257, 65537.
A number of form $2^n+1$ is prime only if it is a Fermat prime.


\Topic{Fermat's Theorem}. Let $m$ be a prime and $x$ and $m$ coprimes, then:
\begin{itemize}
\item $x^{m-1} \equiv 1 \mod m $
\item $x^{k} \mod m = x^{k \mod (m-1)} \mod m$
\item $x^{\phi(m)} \equiv 1 \mod m $
\end{itemize}

%\Topic{Mersenne primes}.  A Mersenne prime is a prime of form $2^n-1$.
%Known Mersenne primes correspond to (necessarily prime) indexes
%2, 3, 5, 7, 13, 17, 19, 31, 61, 89, 107, 127, 521, 607, 1279,
%2203, 2281, 3217, 4253, 4423, 9689, 9941, 11213, 19937,
%21701, 23209, 44497, 86243, 110503, 132049, 216091,
%756839, 859433, 1257787, 1398269, 2976221, 3021377,
%6972593, 13466917.

\Topic{Perfect numbers}.  $n>1$ is called perfect if it equals
sum of its proper divisors and $1$.  Even $n$ is perfect iff $n = 2^{p-1} (2^p - 1)$
and $2^p - 1$ is prime (Mersenne's). No odd perfect numbers are yet found.

\Topic{Carmichael numbers}.
A positive composite $n$ is a Carmichael number
($a^{n-1} \equiv 1 \pmod{n}$ for all $a$: $\gcd(a,n)=1$),
iff $n$ is square-free, and for all prime divisors $p$ of $n$, $p-1$ divides $n-1$.

\Topic{Number/sum of divisors}.
$\tau(p_1^{a_1} \dots p_k^{a_k}) = \prod_{j=1}^k (a_j+1)$. \quad
$\sigma(p_1^{a_1} \dots p_k^{a_k}) = \prod_{j=1}^k \frac{p_j^{a_j+1}-1}{p_j-1}$.

\Topic{Product of divisors}.
$\mu(n) = n^{\frac{\tau(n)}{2}}$
\begin{itemize}
\item if $p$ is a prime, then: $\mu(p^k) = p^{\frac{k(k+1)}{2}}$
\item if $a$ and $b$ are coprimes, then: $\mu(ab) = \mu(a)^{\tau(b)}\mu(b)^{\tau(a)}$
\end{itemize}

\Topic{Euler's phi function}.  $\phi(n)=|\{m \in {\mathbb N}, m \le n, \gcd(m, n) = 1 \}|$.
  
\begin{itemize}
\item   $\phi(mn) = \frac{\phi(m) \phi(n) \gcd(m,n)}{\phi(\gcd(m,n))}$.
\item	$\phi(p) = p - 1$ si $p$ es primo
\item	$\phi(p^a) = p^a(1 - \frac{1}{p}) = p^{a-1} (p-1)$
\item	$\phi(n) = n(1 - \frac{1}{p_1})(1 - \frac{1}{p_2})...(1 - \frac{1}{p_k})$ donde $p_i$ es primo y divide a $n$
\end{itemize}


%\Topic{Fermat's theorem}.  $a^p \equiv a \pmod{p}$ if $p$ is prime. \\
%For any polynomial $f(x)$ with integer coefficients and prime $p$,
%$f(x)^{p^n} \equiv f(x^{p^n}) \pmod{p}$
\Topic{Euler's theorem}. $a^{\phi(n)} \equiv 1\pmod{n}$, if $\gcd(a,n)=1$. \\
\Topic{Wilson's theorem}. $p$ is prime iff $(p - 1)! \equiv -1 \pmod p$.

\Topic{Mobius function}.
$\mu(1) = 1$. $\mu(n) = 0$, if $n$ is not squarefree.
$\mu(n) = (-1)^s$, if $n$ is the product of $s$ distinct primes.
Let $f$, $F$ be functions on positive integers.
If for all $n \in N$, $F(n)=\sum_{d|n} f(d)$, then $f(n) = \sum_{d|n} \mu(d) F(\frac{n}{d})$,
and vice versa. \quad
$\phi(n) = \sum_{d|n} \mu(d) \frac{n}{d}$.
\quad $\sum_{d|n} \mu(d) = 1$. \\
If $f$ is multiplicative, then $\sum_{d|n} \mu(d) f(d) = \prod_{p|n}(1-f(p))$,
$\sum_{d|n} \mu(d)^2 f(d) = \prod_{p|n} (1+f(p))$. \\
$\sum_{d|n} \mu(d) = e(n) = [n==1]$. \\
$S_{f}(n) = \prod_{p = 1} (1+f(p_i) + f(p_i^2 ) +... + f(p_i^{e_i})) $, p - primes(n). 


\Topic{Legendre symbol}. If $p$ is an odd prime, $a \in {\mathbb Z}$, then
$\left(\frac{a}{p}\right)$ equals $0$, if $p | a$; $1$ if $a$ is a quadratic
residue modulo $p$; and $-1$ otherwise.
Euler's criterion:
$\left(\frac{a}{p}\right)=a^{\left(\frac{p-1}{2}\right)} \pmod p$. \\
%$\left(\frac{a}{p}\right) \left(\frac{b}{p}\right) = \left(\frac{ab}{p}\right)$
%Law of Quadratic Reciprocity: for any distinct odd primes $p$ and $q$,
%$\left(\frac{p}{q}\right) \left(\frac{q}{p}\right) = (-1)^{\frac{p-1}{2} \cdot \frac{q-1}{2}}$
\Topic{Jacobi symbol}.  %Generalization of Legendre's symbol to any odd modulus. \\
If $n=p_1^{a_1} \cdots p_k^{a_k}$ is odd, then
$\left(\frac{a}{n}\right) = \prod_{i=1}^k \left(\frac{a}{p_i}\right)^{k_i}$.

%\Topic{Kronecker symbol}.
%Let $a$ be a positive integer, which is not a perfect square and
%$a \equiv 0$ or $1 {\pmod 4}$. \\
%$\left(\frac{a}{2}\right) = \{ 1$, if $a \equiv 1 {\pmod 8}$;
%$-1$, if $a \equiv 5 {\pmod 8} \}$. \\
%$\left(\frac{a}{n}\right) = \prod_{j=1}^k p_j^{k_j}$,
%if gcd$(a,n) \ne 1$ and $n=\prod p_i^{k_i}$.
%$\left(\frac{a}{n}\right)$ equals Jacobi symbol otherwise.

\Topic{Primitive roots}.  If the order of $g$ modulo $m$ (min $n>0$:
$g^n \equiv 1 \pmod{m}$) is $\phi(m)$, then $g$ is called a primitive root.
If $Z_m$ has a primitive root, then it has $\phi(\phi(m))$ distinct primitive
roots. $Z_m$ has a primitive root iff $m$ is one of $2$, $4$,
$p^k$, $2p^k$, where $p$ is an odd prime.
If $Z_m$ has a primitive root $g$, then for all $a$ coprime to $m$,
there exists unique integer $i=\text{ind}_g(a)$ modulo $\phi(m)$,
such that $g^i \equiv a \pmod{m}$.
$\text{ind}_g(a)$ has logarithm-like properties:
$\text{ind}(1) = 0$, $\text{ind}(ab) = \text{ind}(a) + \text{ind}(b)$.

If $p$ is prime and $a$ is not divisible by $p$, then congruence
$x^n \equiv a \pmod{p}$ has $\gcd(n, p-1)$ solutions if
$a^{(p-1)/\gcd(n,p-1)} \equiv 1 \pmod{p}$, and no solutions otherwise.
(Proof sketch: let $g$ be a primitive root, and
$g^i \equiv a \pmod{p}$, $g^u \equiv x \pmod{p}$.
$x^n \equiv a \pmod{p}$ iff $g^{nu} \equiv g^i \pmod{p}$ iff $nu \equiv i \pmod{p}$.)

\Topic{Discrete logarithm problem}.  Find $x$ from $a^x \equiv b \pmod{m}$.
Can be solved in $O(\sqrt{m})$ time and space with a meet-in-the-middle trick.
Let $n = \lceil \sqrt{m} \rceil$, and $x = ny - z$.
Equation becomes $a^{ny} \equiv b a^z \pmod{m}$.  Precompute all values that
the RHS can take for $z = 0, 1, \dots, n-1$, and brute force $y$ on the LHS,
each time checking whether there's a corresponding value for RHS.

\Topic{Pythagorean triples}.  Integer solutions of $x^2 + y^2 = z^2$
All relatively prime triples are given by:
$x=2mn, y=m^2-n^2, z=m^2+n^2$ where $m>n, \gcd(m,n)=1$ and $m \not\equiv n \pmod{2}$.
All other triples are multiples of these.
Equation $x^2 + y^2 = 2z^2$ is equivalent to $(\frac{x+y}{2})^2 + (\frac{x-y}{2})^2 = z^2$.
\begin{itemize}
\item Given an arbitrary pair of integers $m$ and $n$ with $m > n > 0$:\\
$a = m^2 - n^2$, $b = 2mn$, $c = m^2 + n^2$
\item The triple generated by Euclid's formula is primitive if and only if $m$ and $n$ are coprime and not both odd.
\item To generate all Pythagorean triples uniquely:\\
$a = k (m^2 - n^2)$, $b = k(2mn)$, $c = k(m^2 + n^2)$
\item If $m$ and $n$ are two odd integer such that $m > n$, then:\\
$a = mn$, $b = \frac{m^2 - n^2}{2}$, $c = \frac{m^2 + n^2}{2}$
\item If $n=1$ or $2$ there are no solutions. Otherwise\\
$n$ is even: ($(\frac{n^2}{4} - 1)^2 + n^2 = (\frac{n^2}{4} + 1)^2$)\\
$n$ is odd: ($(\frac{n^2 - 1}{2})^2 + n^2 = (\frac{n^2 + 1}{2})^2$)
\end{itemize}


\Topic{Postage stamps/McNuggets problem}.  Let $a$, $b$ be relatively-prime integers.
There are exactly $\frac{1}{2}(a-1)(b-1)$ numbers \emph{not} of form $ax+by$ ($x,y \ge 0$),
and the largest is $(a-1)(b-1)-1 = ab - a - b$.

\Topic{Fermat's two-squares theorem}.  Odd prime $p$ can be represented
as a sum of two squares iff $p \equiv 1 {\pmod 4}$.
A product of two sums of two squares is a sum of two squares.
Thus, $n$ is a sum of two squares iff every prime of
form $p=4k+3$ occurs an even number of times in $n$'s factorization.

\Topic{RSA}. Let $p$ and $q$ be random distinct large primes, $n = pq$.
Choose a small odd integer $e$, relatively prime to $\phi(n) = (p-1)(q-1)$,
and let $d = e^{-1} \pmod{\phi(n)}$. Pairs $(e,n)$ and $(d,n)$ are
the public and secret keys, respectively.
Encryption is done by raising a message $M \in Z_n$ to the power $e$ or $d$,
modulo $n$.